\section{Industry Evidence: Deploying \textsc{SynthTrace} in Practice}

\label{sec:industry-evidence}

This section translates our experimental findings into concrete guidance for practitioners considering \textsc{SynthTrace} in real-world ML pipelines. Rather than positioning synthetic traces as a universal replacement for production data, we describe when and how \textsc{SynthTrace} can be safely and effectively integrated, along with operational guardrails informed by our evaluation.

\subsection{Integration Scenarios}

We identify three practical integration scenarios where \textsc{SynthTrace} provides immediate value.

\paragraph{Augmenting Training Data for ML-Based Observability.}
A common industrial use case is training anomaly detectors or performance models when real kernel traces are scarce or expensive to collect. In this setting, teams can train a diffusion model on a limited set of real traces, generate additional synthetic traces, apply constraint-guided repair, and combine the repaired synthetic data with real traces during downstream training. Our results show that for structured, compute-heavy workloads, this approach achieves near-real performance while doubling effective dataset size, reducing the need for prolonged production tracing.

\paragraph{Validating Trace Consumers and Analysis Pipelines.}
Many organizations maintain complex trace consumers—parsers, aggregators, and analysis pipelines—that are difficult to test comprehensively. Synthetic traces generated by \textsc{SynthTrace} can be used to stress-test these components under diverse but realistic execution patterns. High Top-$K$ accuracy and strong performance on common events indicate that generated traces preserve typical execution behavior while offering broader coverage than limited real datasets.

\paragraph{Testing, Fuzzing, and Development Workflows.}
For non-safety-critical workflows such as testing, fuzzing, and tool prototyping, exact fidelity is often less important than diversity and plausibility. Synthetic traces enable teams to explore rare execution paths, test robustness against unusual sequences, and develop tooling without exposing sensitive production data. In these contexts, even workloads that show macro-F1 degradation remain valuable due to high accuracy on common events and strong Top-$K$ behavior.

\subsection{Cost–Quality Decision Guide}

Our evaluation highlights that context length is the primary driver of synthetic trace quality, but it also introduces significant computational cost. 
In our experimental setup, increasing context length from $L=256$ to $L=4096$ required approximately 16$\times$ more GPU memory and 4$\times$ longer training time. While the quality gains are substantial, they become increasingly expensive beyond $L=4096$, making this configuration a practical upper bound under current hardware constraints.

A cost-effective strategy is therefore to begin with a simple, low-cost configuration and scale context length only when quality requirements justify the additional resource investment.

\subsection{Operational Guardrails (Do / Don’t)}

Based on our results, we summarize clear operational guidance for practitioners.

\begin{center}
\fbox{
\parbox{0.9\columnwidth}{
\textbf{Do}
\begin{itemize}[leftmargin=*,nosep]
\item Use \textsc{SynthTrace} to augment data for deterministic, compute-heavy workloads.
\item Always apply constraint-guided repair before downstream use.
\item Start with minimal feature sets (event + timing) and scale complexity only if needed.
\item Use synthetic traces for testing, development, and privacy-preserving sharing.
\end{itemize}

\vspace{0.5em}

\textbf{Don’t}
\begin{itemize}[leftmargin=*,nosep]
\item Rely on synthetic traces as the sole data source for I/O-heavy or highly nondeterministic systems.
\item Use synthetic data for safety-critical anomaly detection involving rare I/O failures.
\item Increase context length beyond practical limits without considering cost–benefit trade-offs.
\end{itemize}
}
}
\end{center}

These guardrails reflect an important takeaway of our study: synthetic kernel traces are most effective when used as a complement to real data rather than a drop-in replacement.

\subsection{Takeaway for Industry}

From an industrial perspective, \textsc{SynthTrace} offers a pragmatic tool for expanding limited kernel trace datasets, improving testing coverage, and enabling privacy-conscious development workflows. Its strengths lie in structured workloads and long-range temporal modeling, while its limitations highlight open challenges in capturing asynchronous I/O behavior. By following the integration patterns and guardrails outlined above, practitioners can safely leverage synthetic traces without overestimating their fidelity.